% ---- score2 standard preamble (from: python3 score2_tools.py preamble) ----
\documentclass[10pt]{extarticle}
\usepackage[a4paper,margin=1.5cm]{geometry}
\usepackage[T1]{fontenc}
\usepackage{amsmath,amssymb}
\usepackage{graphicx}
\usepackage[export]{adjustbox}
\usepackage{multicol}
\usepackage{enumitem}
\usepackage{booktabs,array,tabularx}
\setlength{\columnsep}{0.7cm}
\setlength{\parindent}{0pt}
\setlength{\parskip}{3pt}
\pagestyle{empty}
\makeatletter
% figure registry (figdefs.tex is generated by the extract step)
\newcommand{\qfigdef}[4]{\@namedef{qf@file@#1}{#2}\@namedef{qf@w@#1}{#3}\@namedef{qf@wide@#1}{#4}}
\newcommand{\qf@check}[1]{\@ifundefined{qf@file@#1}{\PackageError{score2}{Unknown figure id #1}{}}{}}
\newcommand{\qf@img}[4][]{\includegraphics[width=#2\dimexpr\@nameuse{qf@w@#4}pt\relax,max width=#3,max height=0.45\textheight,#1]{\@nameuse{qf@file@#4}}}
% Full-width questions: finalize wraps every question that contains a wide figure in \qfullbegin...\qfullend,
% so the WHOLE question (stem, figures, options) is set across the page in reading order, then the two columns
% resume. The question's head - its number, stem and first figure (\qfullsplit marks the end of the head) -
% is kept on one page, so a number is never left at a page foot with its figure overleaf; the rest may break
% across pages normally, so no big gaps are left. Without a \qfullsplit, the whole question is kept together
% if it is under 45% of a page. Never write these yourself.
\newif\ifqf@full
\newif\ifqf@boxed
\newsavebox{\qf@box}
\newcommand{\qf@open}{\global\qf@fulltrue\par\global\qf@boxedtrue
  \setbox\qf@box\vbox\bgroup\hsize\textwidth\linewidth\textwidth}
\newcommand{\qf@emit}[1]{\par\egroup\global\qf@boxedfalse
  \ifdim\dimexpr\ht\qf@box+\dp\qf@box\relax<#1\textheight\noindent\usebox{\qf@box}\par\else\unvbox\qf@box\fi}
\newcommand{\qfullsplit}{\ifqf@boxed\qf@emit{0.85}\fi}
\newcommand{\qf@close}{\ifqf@boxed\qf@emit{0.45}\fi\global\qf@fullfalse}
\newcommand{\qfullbegin}{\end{multicols}\qf@open}
\newcommand{\qfullend}{\qf@close\begin{multicols}{2}}
\newcommand{\qfullstart}{\qf@open}
\newcommand{\qfulldone}{\qf@close}
% \qfig[scale]{id}: block figure, centred, directly where it is written (size set by the tool).
% Fallback if a wide figure is not inside a full-width question: break out of the columns in place - no floats.
\newcommand{\qfig}[2][1]{\qf@check{#2}%
  \ifqf@full\qf@blk{#1}{#2}\else
  \ifnum\@nameuse{qf@wide@#2}=1 \expandafter\@firstoftwo\else\expandafter\@secondoftwo\fi
  {\end{multicols}\begin{center}\qf@img{#1}{\textwidth}{#2}\end{center}\begin{multicols}{2}}%
  {\qf@blk{#1}{#2}}\fi}
\newcommand{\qf@blk}[2]{\par\smallskip{\centering\qf@img{#1}{\linewidth}{#2}\par}\smallskip}
% \qfigi[scale]{id}: inline figure (inside option lists / tables / minipages); never breaks columns.
\newcommand{\qfigi}[2][1]{\qf@check{#2}\qf@img[valign=c]{#1}{\linewidth}{#2}}
\makeatother
% structure macros - lint reads these to check completeness and figure placement
\newcommand{\papersec}[1]{\par\bigskip{\centering\Large\bfseries PAPER-#1\par}\medskip}
\newcommand{\partsec}[2]{\par\medskip\noindent{\large\bfseries PART-#1: #2}\par\smallskip}
\newcommand{\sect}[1]{\par\smallskip\noindent\textbf{#1}\par}
\newcommand{\qn}[1]{\par\smallskip\noindent\textbf{#1.}~}
\newcommand{\paraq}[2]{\par\smallskip\noindent\textbf{Paragraph for Questions #1 and #2}\par}
\newenvironment{qoptions}{\begin{itemize}[leftmargin=1.8em,itemsep=1pt,topsep=2pt,label={}]}{\end{itemize}}
\newcommand{\opt}[2]{\item[(#1)] #2}
% figure registry (generated by score2_tools.py finalize)
\qfigdef{p1_phy_s1_q01_1}{Score-II - Test 21 (14 Mar 2024)-img/p1_phy_s1_q01_1.png}{58.3}{0}
\qfigdef{p1_phy_s1_q01_2}{Score-II - Test 21 (14 Mar 2024)-img/p1_phy_s1_q01_2.png}{58.3}{0}
\qfigdef{p1_phy_s1_q02_1}{Score-II - Test 21 (14 Mar 2024)-img/p1_phy_s1_q02_1.png}{121.8}{0}
\qfigdef{p1_phy_s1_q03_1}{Score-II - Test 21 (14 Mar 2024)-img/p1_phy_s1_q03_1.png}{82.0}{0}
\qfigdef{p1_phy_s1_q04_1}{Score-II - Test 21 (14 Mar 2024)-img/p1_phy_s1_q04_1.png}{87.2}{0}
\qfigdef{p1_phy_s1_q05_1}{Score-II - Test 21 (14 Mar 2024)-img/p1_phy_s1_q05_1.png}{96.2}{0}
\qfigdef{p1_phy_s2_q07_1}{Score-II - Test 21 (14 Mar 2024)-img/p1_phy_s2_q07_1.png}{134.0}{0}
\qfigdef{p1_phy_s2_q07_2}{Score-II - Test 21 (14 Mar 2024)-img/p1_phy_s2_q07_2.png}{69.7}{0}
\qfigdef{p1_phy_s2_q07_3}{Score-II - Test 21 (14 Mar 2024)-img/p1_phy_s2_q07_3.png}{71.0}{0}
\qfigdef{p1_phy_s2_q07_4}{Score-II - Test 21 (14 Mar 2024)-img/p1_phy_s2_q07_4.png}{123.8}{0}
\qfigdef{p1_phy_s2_q07_5}{Score-II - Test 21 (14 Mar 2024)-img/p1_phy_s2_q07_5.png}{78.6}{0}
\qfigdef{p1_phy_s2_q10_1}{Score-II - Test 21 (14 Mar 2024)-img/p1_phy_s2_q10_1.png}{104.8}{0}
\qfigdef{p1_phy_s2_q10_2}{Score-II - Test 21 (14 Mar 2024)-img/p1_phy_s2_q10_2.png}{105.6}{0}
\qfigdef{p1_phy_s2_q10_3}{Score-II - Test 21 (14 Mar 2024)-img/p1_phy_s2_q10_3.png}{104.3}{0}
\qfigdef{p1_phy_s2_q10_4}{Score-II - Test 21 (14 Mar 2024)-img/p1_phy_s2_q10_4.png}{103.9}{0}
\qfigdef{p1_phy_s3_q03_1}{Score-II - Test 21 (14 Mar 2024)-img/p1_phy_s3_q03_1.png}{153.7}{0}
\qfigdef{p1_phy_s3_q04_1}{Score-II - Test 21 (14 Mar 2024)-img/p1_phy_s3_q04_1.png}{197.7}{0}
\qfigdef{p1_phy_s3_q06_1}{Score-II - Test 21 (14 Mar 2024)-img/p1_phy_s3_q06_1.png}{129.4}{0}
\qfigdef{p1_chem_s1_q05_1}{Score-II - Test 21 (14 Mar 2024)-img/p1_chem_s1_q05_1.png}{325.8}{1}
\qfigdef{p1_chem_s1_q06_1}{Score-II - Test 21 (14 Mar 2024)-img/p1_chem_s1_q06_1.png}{170.7}{0}
\qfigdef{p1_chem_s1_q06_2}{Score-II - Test 21 (14 Mar 2024)-img/p1_chem_s1_q06_2.png}{79.9}{0}
\qfigdef{p1_chem_s1_q06_3}{Score-II - Test 21 (14 Mar 2024)-img/p1_chem_s1_q06_3.png}{245.2}{0}
\qfigdef{p1_chem_s1_q06_4}{Score-II - Test 21 (14 Mar 2024)-img/p1_chem_s1_q06_4.png}{236.8}{0}
\qfigdef{p1_chem_s2_q07_1}{Score-II - Test 21 (14 Mar 2024)-img/p1_chem_s2_q07_1.png}{88.6}{0}
\qfigdef{p1_chem_s2_q07_2}{Score-II - Test 21 (14 Mar 2024)-img/p1_chem_s2_q07_2.png}{91.9}{0}
\qfigdef{p1_chem_s2_q07_3}{Score-II - Test 21 (14 Mar 2024)-img/p1_chem_s2_q07_3.png}{90.5}{0}
\qfigdef{p1_chem_s2_q07_4}{Score-II - Test 21 (14 Mar 2024)-img/p1_chem_s2_q07_4.png}{84.9}{0}
\qfigdef{p1_chem_s2_q07_5}{Score-II - Test 21 (14 Mar 2024)-img/p1_chem_s2_q07_5.png}{85.0}{0}
\qfigdef{p1_chem_s3_q06_1}{Score-II - Test 21 (14 Mar 2024)-img/p1_chem_s3_q06_1.png}{163.2}{0}
\qfigdef{p1_chem_s3_q07_1}{Score-II - Test 21 (14 Mar 2024)-img/p1_chem_s3_q07_1.png}{216.3}{0}
\qfigdef{p1_chem_s3_q08_2}{Score-II - Test 21 (14 Mar 2024)-img/p1_chem_s3_q08_2.png}{331.4}{1}
\qfigdef{p2_phy_s1_q01_1}{Score-II - Test 21 (14 Mar 2024)-img/p2_phy_s1_q01_1.png}{72.3}{0}
\qfigdef{p2_phy_s1_q04_1}{Score-II - Test 21 (14 Mar 2024)-img/p2_phy_s1_q04_1.png}{161.1}{0}
\qfigdef{p2_phy_s2_q05_1}{Score-II - Test 21 (14 Mar 2024)-img/p2_phy_s2_q05_1.png}{75.3}{0}
\qfigdef{p2_phy_s2_q06_2}{Score-II - Test 21 (14 Mar 2024)-img/p2_phy_s2_q06_2.png}{113.9}{0}
\qfigdef{p2_phy_s2_q06_1}{Score-II - Test 21 (14 Mar 2024)-img/p2_phy_s2_q06_1.png}{47.0}{0}
\qfigdef{p2_phy_s2_q09_1}{Score-II - Test 21 (14 Mar 2024)-img/p2_phy_s2_q09_1.png}{165.5}{0}
\qfigdef{p2_phy_s2_q10_1}{Score-II - Test 21 (14 Mar 2024)-img/p2_phy_s2_q10_1.png}{152.6}{0}
\qfigdef{p2_phy_s3_q03_1}{Score-II - Test 21 (14 Mar 2024)-img/p2_phy_s3_q03_1.png}{156.7}{0}
\qfigdef{p2_phy_s3_q05_1}{Score-II - Test 21 (14 Mar 2024)-img/p2_phy_s3_q05_1.png}{153.3}{0}
\qfigdef{p2_phy_s3_q06_1}{Score-II - Test 21 (14 Mar 2024)-img/p2_phy_s3_q06_1.png}{74.9}{0}
\qfigdef{p2_phy_s3_q08_1}{Score-II - Test 21 (14 Mar 2024)-img/p2_phy_s3_q08_1.png}{183.5}{0}
\qfigdef{p2_chem_s1_q04_A}{Score-II - Test 21 (14 Mar 2024)-img/p2_chem_s1_q04_A.png}{210.1}{0}
\qfigdef{p2_chem_s1_q04_B}{Score-II - Test 21 (14 Mar 2024)-img/p2_chem_s1_q04_B.png}{209.2}{0}
\qfigdef{p2_chem_s1_q04_C}{Score-II - Test 21 (14 Mar 2024)-img/p2_chem_s1_q04_C.png}{207.7}{0}
\qfigdef{p2_chem_s1_q04_D}{Score-II - Test 21 (14 Mar 2024)-img/p2_chem_s1_q04_D.png}{204.4}{0}
\qfigdef{p2_chem_s2_q09_1}{Score-II - Test 21 (14 Mar 2024)-img/p2_chem_s2_q09_1.png}{189.8}{0}
\qfigdef{p2_chem_s3_q03_2}{Score-II - Test 21 (14 Mar 2024)-img/p2_chem_s3_q03_2.png}{152.5}{0}
\qfigdef{p2_chem_s3_q06_1}{Score-II - Test 21 (14 Mar 2024)-img/p2_chem_s3_q06_1.png}{168.2}{0}
\qfigdef{p2_chem_s3_q07_2}{Score-II - Test 21 (14 Mar 2024)-img/p2_chem_s3_q07_2.png}{186.4}{0}
\qfigdef{p2_chem_s3_q08_1}{Score-II - Test 21 (14 Mar 2024)-img/p2_chem_s3_q08_1.png}{227.0}{0}
% ---- body skeleton ----
% \begin{document}
% \begin{center}{\Large\bfseries JEE (Main+Advanced) Score-II --- Test NN (DD Mon YYYY)}\end{center}
% \begin{multicols}{2}
%   \papersec{1}   (only for two-paper tests)
%   \partsec{1}{PHYSICS} \sect{Section-I (i): Single Correct} \qn{1} ... \qfig{p1_phy_q01_1} ...
% \end{multicols}
% \end{document}

\begin{document}
\begin{center}{\Large\bfseries JEE (Main+Advanced) Score-II --- Test 21 (14 Mar 2024)}\end{center}
\begin{multicols}{2}

\papersec{1}
\partsec{1}{PHYSICS}
\sect{Section-I (i): One or More Correct}

\qn{1} A medium having dielectric constant $K>1$ fills the space between the plates of a parallel plate capacitor. The plates have large area, and the distance between them is $d$. The capacitor is connected to a battery of voltage $V$, as shown in Figure (a). Now, the left plate is moved by a distance of $\dfrac{d}{2}$ from its original position, as shown in Figure (b).
\qfig{p1_phy_s1_q01_1}
\qfig{p1_phy_s1_q01_2}
In the process of going from the configuration depicted in Figure (a) to that in Figure (b), which of the following statement(s) is(are) correct?
\begin{qoptions}
\opt{A}{The electric field inside the dielectric material is reduced by a factor of $K+2$.}
\opt{B}{The capacitance is decreased by a factor of $\dfrac{2}{K+2}$.}
\opt{C}{The voltage between the capacitor plates remains same.}
\opt{D}{The energy of capacitor decreases.}
\end{qoptions}

\qn{2} The figure shows a circuit having eight resistances of $1\,\Omega$ each, labelled $R_1$ to $R_8$, and two ideal batteries with voltages $\varepsilon_1 = 1$ V and $\varepsilon_2 = 2$ V.
\qfig{p1_phy_s1_q02_1}
Which of the following statement(s) is(are) correct?
\begin{qoptions}
\opt{A}{The magnitude of current flowing through $R_1$ is 0.8 $A$.}
\opt{B}{The magnitude of current flowing through $R_2$ is 0.2 $A$.}
\opt{C}{The magnitude of current flowing through $R_3$ is 1.2 $A$.}
\opt{D}{The magnitude of current flowing through $R_5$ is 0.6 $A$.}
\end{qoptions}

\qn{3} An ideal gas of density $\rho = 0.2\ kg\,m^{-3}$ enters a chimney of height h at the rate of $\alpha = 0.8\ kg\,s^{-1}$ from its lower end, and escapes through the upper end as shown in the figure. The cross-sectional area of the lower end is $A_1 = 0.1\ m^2$ and the upper end is $A_2 = 0.4\ m^2$. The pressure and the temperature of the gas at the lower end are 600 $Pa$ and 300 $K$, respectively, while its temperature at the upper end is 150 $K$. The chimney is heat insulated so that the gas undergoes adiabatic expansion. Take $g=10\ ms^{-2}$ and the ratio of specific heats of the gas $\gamma = \dfrac32$. Ignore atmospheric pressure. Which of the following statement(s) is(are) correct?
\qfig{p1_phy_s1_q03_1}
\begin{qoptions}
\opt{A}{The pressure of the gas at the upper end of the chimney is 75 $Pa$.}
\opt{B}{The velocity of the gas at the both the ends of the chimney is 40 $ms^{-1}$.}
\opt{C}{The height of the chimney is 590 $m$.}
\opt{D}{The density of the gas at the upper end is 0.05 $kg\,m^{-3}$.}
\end{qoptions}

\qn{4} Three plane mirrors form an equilateral triangle with each side of length $L$. There is a small hole at a distance $\ell > 0$ from one of the corners as shown in the figure. A ray of light is passed through the hole at an angle $\theta$ and can only come out through the same hole. The cross section of the mirror configuration and the ray of light lie on the same plane. Which of the following statement(s) is(are) correct?
\qfig{p1_phy_s1_q04_1}
\begin{qoptions}
\opt{A}{The ray of light will come out for $\theta = 30^\circ$, for $0 < \ell < L$.}
\opt{B}{There is an angle for $\ell = \dfrac{L}{2}$ at which the ray of light will come out after two reflections.}
\opt{C}{The ray of light will NEVER come out for $\theta = 60^\circ$, and $\ell = \dfrac{L}{3}$.}
\opt{D}{The ray of light will come out for $\theta = 60^\circ$, and $0 < \ell < \dfrac{L}{2}$ after six reflections.}
\end{qoptions}

\qn{5} Six charges are placed around a regular hexagon of side length $a$ as shown in the figure. Five of them have charge $q$, and the remaining one has charge $x$. The perpendicular from each charge to the nearest hexagon side passes through the center O of the hexagon and is bisected by the side.
\qfig{p1_phy_s1_q05_1}
Which of the following statement(s) is(are) correct in SI units?
\begin{qoptions}
\opt{A}{When $x = q$, the magnitude of the electric field at O is zero.}
\opt{B}{When $x = -q$, the magnitude of the electric field at O is $\dfrac{q}{6\pi\epsilon_0 a^2}$.}
\opt{C}{When $x = 2q$, the potential at O is $\dfrac{7q}{4\sqrt3\pi\epsilon_0 a}$.}
\opt{D}{When $x = -3q$, the potential at O is $\dfrac{3q}{4\sqrt3\pi\epsilon_0 a}$.}
\end{qoptions}

\qn{6} The binding energy of nucleons in a nucleus can be affected by the pairwise Coulomb repulsion. Assume that all nucleons are uniformly distributed inside the nucleus. Let the binding energy of a proton be $E_b^p$ and the binding energy of a neutron be $E_b^n$ in the nucleus. Which of the following statement(s) is(are) correct?
\begin{qoptions}
\opt{A}{$E_b^n - E_b^p$ increases with increases in atomic number.}
\opt{B}{$E_b^n - E_b^p$ is proportional to $A^{-1/3}$ where $A$ is the mass number of the nucleus.}
\opt{C}{$E_b^p - E_b^n$ is positive.}
\opt{D}{$E_b^p$ increases if the nucleus undergoes a beta decay emitting a positron.}
\end{qoptions}

\sect{Section-I (ii): Matching List (Single Correct)}

\qn{7} A small circular loop of area $A$ and resistance $R$ is fixed on a horizontal $xy$-plane with the center of the loop always on the axis $\hat n$ of a long solenoid. The solenoid has $m$ turns per unit length and carries current $I$ counterclockwise as shown in the figure. The magnetic field due to the solenoid is in $\hat n$ direction. List-I gives time dependences of $\hat n$ in terms of a constant angular frequency $\omega$. List-II gives the torques experienced by the circular loop at time $t = \dfrac{\pi}{6\omega}$, Let $\alpha = \dfrac{A^2\mu_0^2 m^2 I^2\omega}{2R}$. Which one of the following options is correct?
\qfig{p1_phy_s2_q07_1}
\end{multicols}
\begin{center}
\begin{tabular}{|c|l|c|l|}
\hline
\multicolumn{2}{|c|}{List-I} & \multicolumn{2}{c|}{List-II} \\
\hline
(I) & $\dfrac{1}{\sqrt2}\left(\sin\omega t\,\hat\jmath + \cos\omega t\,\hat k\right)$ & (P) & $0$ \\
\hline
(II) & $\dfrac{1}{\sqrt2}\left(\sin\omega t\,\hat\imath + \cos\omega t\,\hat\jmath\right)$ & (Q) & $-\dfrac{\alpha}{4}\hat\imath$ \\
\hline
(III) & $\dfrac{1}{\sqrt2}\left(\sin\omega t\,\hat\imath + \cos\omega t\,\hat k\right)$ & (R) & $\dfrac{3\alpha}{4}\hat\imath$ \\
\hline
(IV) & $\dfrac{1}{\sqrt2}\left(\cos\omega t\,\hat\imath + \sin\omega t\,\hat k\right)$ & (S) & $\dfrac{\alpha}{4}\hat\jmath$ \\
\hline
& & (T) & $-\dfrac{3\alpha}{4}\hat\imath$ \\
\hline
\end{tabular}
\end{center}
\begin{multicols}{2}
\begin{qoptions}
\opt{A}{I$\to$Q; II$\to$P; III$\to$S; IV$\to$T}
\opt{B}{I$\to$S; II$\to$T; III$\to$Q; IV$\to$P}
\opt{C}{I$\to$Q; II$\to$P; III$\to$S; IV$\to$R}
\opt{D}{I$\to$T; II$\to$Q; III$\to$P; IV$\to$R}
\end{qoptions}

\qn{8} List I describes four systems, each with two particles $A$ and $B$ in relative motion as shown in figure. List II gives possible magnitudes of their relative velocities (in $ms^{-1}$) at time $t = \dfrac{\pi}{3}$ s. Which one of the following options is correct?
\end{multicols}
\begin{center}
\begin{tabular}{|c|p{7cm}|c|l|}
\hline
\multicolumn{2}{|c|}{List-I} & \multicolumn{2}{c|}{List-II} \\
\hline
(I) & $A$ and $B$ are moving on a horizontal circle of radius 1 $m$ with uniform angular speed $\omega = 1\ rad\,s^{-1}$. The initial angular positions of $A$ and $B$ at time $t=0$ are $\theta=0$ and $\theta = \dfrac{\pi}{2}$ respectively.
\qfigi{p1_phy_s2_q07_2} & (P) & $\dfrac{\sqrt3+1}{2}$ \\
\hline
(II) & Projectiles $A$ and $B$ are fired (in the same vertical plane) at $t=0$ and $t=0.1$ s respectively, with the same speed $v = \dfrac{5\pi}{\sqrt2}\ m\,s^{-1}$ and at $45^\circ$ from the horizontal plane. The initial separation between $A$ and $B$ is large enough so that they do not collide, ($g=10\ m\,s^{-2}$).
\qfigi{p1_phy_s2_q07_3} & (Q) & $\dfrac{\sqrt3-1}{\sqrt2}$ \\
\hline
(III) & Two harmonic oscillators $A$ and $B$ moving in the $x$ direction according to $x_A = x_0\sin\dfrac{t}{t_0}$ and $x_B = x_0\sin\left(\dfrac{t}{t_0}+\dfrac{\pi}{2}\right)$ respectively, starting from $t=0$. Take $x_0=1$ m, $t_0=1$ s.
\qfigi{p1_phy_s2_q07_4} & (R) & $\sqrt{10}$ \\
\hline
(IV) & Particle $A$ is rotating in a horizontal circular path of radius 1 $m$ on the $xy$ plane, with constant angular speed $\omega=1\ rad\,s^{-1}$. Particle $B$ is moving up at a constant speed 3 $m\,s^{-1}$ in the vertical direction as shown in the figure. (Ignore gravity.)
\qfigi{p1_phy_s2_q07_5} & (S) & $\sqrt2$ \\
\hline
& & (T) & $\sqrt{25\pi^2+1}$ \\
\hline
\end{tabular}
\end{center}
\begin{multicols}{2}
\begin{qoptions}
\opt{A}{I$\to$R; II$\to$T; III$\to$P; IV$\to$S}
\opt{B}{I$\to$S; II$\to$P; III$\to$Q; IV$\to$R}
\opt{C}{I$\to$S; II$\to$T; III$\to$P; IV$\to$R}
\opt{D}{I$\to$T; II$\to$P; III$\to$R; IV$\to$S}
\end{qoptions}

\qn{9} List I describes thermodynamic processes in four different systems. List II gives the magnitudes (either exactly or as a close approximation) of possible changes in the internal energy of the system due to the process. Which one of the following options is correct?
\end{multicols}
\begin{center}
\begin{tabular}{|c|p{7.3cm}|c|l|}
\hline
\multicolumn{2}{|c|}{List-I} & \multicolumn{2}{c|}{List-II} \\
\hline
(I) & $10^{-3}$ kg of water at $100^\circ$C is converted to steam at the same temperature, at a pressure of $10^5\ Pa$. The volume of the system changes from $10^{-6}\ m^3$ to $10^{-3}\ m^3$ in the process. Latent heat of water = 2250 kJ/kg. & (P) & 2 $kJ$ \\
\hline
(II) & 0.2 moles of a rigid diatomic ideal gas with volume $V$ at temperature 500 $K$ undergoes an isobaric expansion to volume 3 $V$. Assume $R = 8.0\ J\,mol^{-1}K^{-1}$. & (Q) & 7 $kJ$ \\
\hline
(III) & One mole of a monatomic ideal gas is compressed adiabatically from volume $V = \dfrac13\ m^3$ and pressure 2 $kPa$ to volume $\dfrac{V}{8}$ & (R) & 4 $kJ$ \\
\hline
(IV) & Three moles of a diatomic ideal gas whose molecules can vibrate, is given 9 $kJ$ of heat and undergoes isobaric expansion. & (S) & 5 $kJ$ \\
\hline
& & (T) & 3 $kJ$ \\
\hline
\end{tabular}
\end{center}
\begin{multicols}{2}
\begin{qoptions}
\opt{A}{I$\to$T; II$\to$R; III$\to$S; IV$\to$Q}
\opt{B}{I$\to$S; II$\to$P; III$\to$T; IV$\to$P}
\opt{C}{I$\to$P; II$\to$R; III$\to$T; IV$\to$Q}
\opt{D}{I$\to$Q; II$\to$R; III$\to$S; IV$\to$T}
\end{qoptions}

\qn{10} List I contains four combinations of two lenses (1 and 2) whose focal lengths (in cm) are indicated in the figures. In all cases, the object is placed 20 $cm$ from the first lens on the left, and the distance between the two lenses is 5 $cm$. List II contains the positions of the final images. Which one of the following options is correct?
\end{multicols}
\begin{center}
\begin{tabular}{|c|c|c|p{4.3cm}|}
\hline
\multicolumn{2}{|c|}{List-I} & \multicolumn{2}{c|}{List-II} \\
\hline
(I) & \qfigi{p1_phy_s2_q10_1} & (P) & Final image is formed at 7.5 $cm$ on the right side of lens 2. \\
\hline
(II) & \qfigi{p1_phy_s2_q10_2} & (Q) & Final image is formed at 60.0 $cm$ on the right side of lens 2. \\
\hline
(III) & \qfigi{p1_phy_s2_q10_3} & (R) & Final image is formed at 30.0 $cm$ on the left side of lens 2. \\
\hline
(IV) & \qfigi{p1_phy_s2_q10_4} & (S) & Final image is formed at 6.0 $cm$ on the right side of lens 2. \\
\hline
& & (T) & Final image is formed at 30.0 $cm$ on the right side of lens 2. \\
\hline
\end{tabular}
\end{center}
\begin{multicols}{2}
\begin{qoptions}
\opt{A}{I$\to$P; II$\to$R; III$\to$Q; IV$\to$T}
\opt{B}{I$\to$Q; II$\to$P; III$\to$T; IV$\to$S}
\opt{C}{I$\to$P; II$\to$T; III$\to$R; IV$\to$Q}
\opt{D}{I$\to$T; II$\to$S; III$\to$Q; IV$\to$R}
\end{qoptions}

\sect{Section-II: Numerical Value}

\qn{1} Two spherical stars $A$ and $B$ have densities $\rho_A$ and $\rho_B$, respectively. $A$ and $B$ have the same radius, and their masses $M_A$ and $M_B$ are related by $M_B = 8M_A$. Due to an interaction process, star $A$ gains some mass, so that its radius is doubled, while its spherical shape is retained, and its density remains $\rho_A$. The entire mass gain by $A$ is from an outer spherical shell of $B$. If $v_A$ and $v_B$ are the escape velocities from $A$ and $B$ after the interaction process, the ratio $\dfrac{v_B}{v_A} = n$. The value of $n$ is\underline{\hspace{2em}}

\qn{2} The minimum kinetic energy needed by an alpha particle to cause the nuclear reaction ${}^{16}_7$N $+ {}^4_2$He $\to {}^1_1$H $+ {}^{19}_8$O in a laboratory frame is $n$ (in MeV). Assume ${}^{16}_7$N is at rest in the laboratory frame. The masses of ${}^{16}_7$N, ${}^4_2$He, ${}^1_1$H and ${}^{19}_8$O can be taken to be 16.006 u, 4.003 u, 1.008 u and 19.003 u, respectively, where 1 u $= 930\ MeVc^{-2}$. The value of $n$ is\underline{\hspace{2em}}

\qn{3} In the following circuit $C_1 = 12\ \mu F$, $C_2 = C_3 = 4\ \mu F$ and $C_4 = C_5 = 2\ \mu F$. The energy stored in $C_1$ is\underline{\hspace{2em}} $\mu J$.
\qfig{p1_phy_s3_q03_1}

\qn{4} A rod of length 2 $cm$ makes an angle $\dfrac{2\pi}{3}$ rad with the principal axis of a thin convex lens. The lens has a focal length of 10 $cm$ and is placed at a distance of 15 cm from the object as shown in the figure. The angle made by it with respect to the principal axis is $\alpha$ rad. The value of $\tan\alpha$ is\underline{\hspace{2em}}.
\qfig{p1_phy_s3_q04_1}

\qn{5} At time $t=0$, a disk of radius $\sqrt\pi\,m$ starts to roll without slipping on a horizontal plane with a constant angular acceleration of $\alpha = \dfrac23\ rad\,s^{-2}$. A small stone of mass $\dfrac3\pi$ kg is stuck to the disk. At $t=0$, it is at the contact point of the disk and the plane. Later, at time $t=\sqrt\pi$ s, the stone detaches itself and flies off tangentially from the disk. The velocity of stone at that moment is\underline{\hspace{2em}} m/s. [Take $g=10\ m\,s^{-2}$.]

\qn{6} A solid sphere of mass 1 $kg$ and radius 1 $m$ rolls without slipping on a fixed inclined plane with an angle of inclination $\theta = 30^\circ$ from the horizontal. Two forces of magnitude 2 $N$ each, parallel to the incline, act on the sphere, both at distance $r = 0.5$ $m$ from the center of the sphere, as shown in the figure. The acceleration of the sphere down the plane is\underline{\hspace{2em}} $ms^{-2}$. (Take $g=10\ m\,s^{-2}$.)
\qfig{p1_phy_s3_q06_1}

\qn{7} Consider an LC circuit, with inductance $L=0.2$ H and capacitance $C=10^{-3}$ F, kept on a plane. The area of the circuit is 1 $m^2$. It is placed in a constant magnetic field of strength $B_0$ which is perpendicular to the plane of the circuit. At time $t=0$, the magnetic field strength starts increasing linearly as $B = B_0+\beta t$ with $\beta = 0.04\ Ts^{-1}$. The maximum magnitude of the current in the circuit is\underline{\hspace{2em}} mA.

\qn{8} A projectile is fired from horizontal ground with speed v and projection angle $\theta$. When the acceleration due to gravity is g, the time of flight of the projectile is t. If at the highest point in its trajectory, the projectile enters a different region where the effective acceleration due to gravity is $g' = \dfrac{g}{0.8}$, then the new time of flight is $t' = nt$. The value of $n$ is\underline{\hspace{2em}}.

\partsec{2}{CHEMISTRY}
\sect{Section-I (i): One or More Correct}

\qn{1} For diatomic molecules, the correct statement(s) about the molecular orbitals formed by the overlap to two $2p_z$ orbitals is(are)
\begin{qoptions}
\opt{A}{$\sigma$ orbital has a total of two nodal planes.}
\opt{B}{$\pi^*$ orbital has one node in the $xy$-plane containing the molecular axis.}
\opt{C}{$\sigma^*$ orbital has one node in the $xz$-plane containing the molecular axis.}
\opt{D}{$\pi$ orbital has one node in the plane which is perpendicular to the molecular axis and goes through the center of the molecule.}
\end{qoptions}

\qn{2} The correct option(s) related to adsorption processes is(are)
\begin{qoptions}
\opt{A}{Chemisorption results in a unimolecular layer.}
\opt{B}{The enthalpy change during chemisorption is in the range of 100 to 140 $kJ\,mol^{-1}$.}
\opt{C}{Chemisorption is an endothermic process.}
\opt{D}{Lowering the temperature favors physisorption processes.}
\end{qoptions}

\qn{3} The electrochemical extraction of aluminum from bauxite ore involves.
\begin{qoptions}
\opt{A}{the neutralization of aluminate solution by passing CO$_2$ gas to precipitate hydrated alumina (Al$_2$O$_3$.3H$_2$O)}
\opt{B}{the dissolution of Al$_2$O$_3$ in hot aqueous NaOH.}
\opt{C}{the electrolysis of Al$_2$O$_3$ mixed with Na$_3$AlF$_6$ to give Al and CO$_2$.}
\opt{D}{the reaction of Al$_2$O$_3$ with coke (C) at a temperature $>2500^\circ$C.}
\end{qoptions}

\qn{4} The treatment of galena with HNO$_3$ produces a gas that is
\begin{qoptions}
\opt{A}{paramagnetic}
\opt{B}{linear in geometry}
\opt{C}{a neutral oxide}
\opt{D}{colorless}
\end{qoptions}

\qfullbegin
\qn{5} Considering the reaction sequence given below, the correct statement(s) is(are)
\qfig{p1_chem_s1_q05_1}
\qfullsplit
\begin{qoptions}
\opt{A}{$\mathbf{P}$ can be reduced to a primary alcohol using LiAlH$_4$.}
\opt{B}{Treating $\mathbf{P}$ with conc. NH$_4$OH solution followed by acidification gives $\mathbf{Q}$.}
\opt{C}{Treating $\mathbf{Q}$ with a solution of NaNO$_2$ in aq. HCl liberates N$_2$.}
\opt{D}{$\mathbf{P}$ is more acidic than CH$_3$CH$_2$COOH.}
\end{qoptions}
\qfullend

\qn{6} Consider the following reaction sequence,
\qfig{p1_chem_s1_q06_1}
the correct option(s) is(are)
\begin{qoptions}
\opt{A}{$\mathbf{P}$ = H$_2$/Pd, ethanol \hspace{2em} $\mathbf{R}$ = NaNO$_2$/HCl \hspace{2em} $\mathbf{U}$ = 1. H$_3$PO$_2$\newline\hbox{}\hspace{10.2em}2. KMnO$_4$ - KOH, heat}
\opt{B}{$\mathbf{P}$ = Sn/HCl \hspace{2em} $\mathbf{R}$ = HNO$_2$ \hspace{2em} $\mathbf{S}$ =\\ \qfigi{p1_chem_s1_q06_2}}
\opt{C}{$\mathbf{S}$ =\\ \qfigi[0.85]{p1_chem_s1_q06_3}\\ $\mathbf{U}$ = 1. CH$_3$CH$_2$OH\\ \hbox{}\hspace{1.5em}2. KMnO$_4$ - KOH, heat}
\opt{D}{$\mathbf{Q}$ =\\ \qfigi[0.85]{p1_chem_s1_q06_4}}
\end{qoptions}

\sect{Section-I (ii): Matching List (Single Correct)}

\qn{7} Match the rate expressions in LIST-I for the decomposition of X with the corresponding profiles provided in LIST-II. X$_s$ and k constants having appropriate units.
\end{multicols}
\begin{center}
\begin{tabular}{|c|p{5.5cm}|c|c|}
\hline
\multicolumn{2}{|c|}{LIST-I} & \multicolumn{2}{c|}{LIST-II} \\
\hline
(I) & rate $= \dfrac{k[X]}{X_s+[X]}$ under all possible initial concentration of X & (P) & \qfigi{p1_chem_s2_q07_1} \\
\hline
(II) & rate $= \dfrac{k[X]}{X_s+[X]}$ where initial concentration of X are much less than $X_s$ & (Q) & \qfigi{p1_chem_s2_q07_2} \\
\hline
(III) & rate $= \dfrac{k[X]}{X_s+[X]}$ where initial concentration of X are much higher than $X_s$ & (R) & \qfigi{p1_chem_s2_q07_3} \\
\hline
(IV) & rate $= \dfrac{k[X]^2}{X_s+[X]}$ where initial concentration of X is much higher than $X_s$ & (S) & \qfigi{p1_chem_s2_q07_4} \\
\hline
& & (T) & \qfigi{p1_chem_s2_q07_5} \\
\hline
\end{tabular}
\end{center}
\begin{multicols}{2}
\begin{qoptions}
\opt{A}{I$\to$R; II$\to$S; III$\to$S; IV$\to$T}
\opt{B}{I$\to$P; II$\to$Q; III$\to$Q; IV$\to$R}
\opt{C}{I$\to$R; II$\to$S; III$\to$Q; IV$\to$R}
\opt{D}{I$\to$R; II$\to$Q; III$\to$S; IV$\to$T}
\end{qoptions}

\qn{8} LIST-I contains compounds and LIST-II contains reaction
\end{multicols}
\begin{center}
\begin{tabular}{|c|c|c|l|}
\hline
\multicolumn{2}{|c|}{LIST-I} & \multicolumn{2}{c|}{LIST-II} \\
\hline
(I) & H$_2$O$_2$ & (P) & Mg(HCO$_3$)$_2$ + Ca(OH)$_2$ $\longrightarrow$ \\
\hline
(II) & Mg(OH)$_2$ & (Q) & BaO$_2$ + H$_2$SO$_4$ $\longrightarrow$ \\
\hline
(III) & BaCl$_2$ & (R) & Ca(OH)$_2$ + MgCl$_2$ $\longrightarrow$ \\
\hline
(IV) & CaCO$_3$ & (S) & BaO$_2$ + HCl $\longrightarrow$ \\
\hline
& & (T) & Ca(HCO$_3$)$_2$ + Ca(OH)$_2$ $\longrightarrow$ \\
\hline
\end{tabular}
\end{center}
\begin{multicols}{2}
Match each compound in LIST-I with its formation reaction(s) in LIST-II, and choose the correct option
\begin{qoptions}
\opt{A}{I$\to$Q; II$\to$P; III$\to$S; IV$\to$R}
\opt{B}{I$\to$T; II$\to$P; III$\to$Q; IV$\to$R}
\opt{C}{I$\to$Q; II$\to$R; III$\to$S; IV$\to$P}
\opt{D}{I$\to$T; II$\to$R; III$\to$Q; IV$\to$P}
\end{qoptions}

\qn{9} LIST-I contains metal species and LIST-II contains their properties.
\end{multicols}
\begin{center}
\begin{tabular}{|c|c|c|l|}
\hline
\multicolumn{2}{|c|}{LIST-I} & \multicolumn{2}{c|}{LIST-II} \\
\hline
(I) & [Cr(CN)$_6$]$^{4-}$ & (P) & $t_{2g}$ orbitals contain 4 electrons \\
\hline
(II) & [RuCl$_6$]$^{2-}$ & (Q) & $\mu$(spin-only) $=4.9$ BM \\
\hline
(III) & [Cr(H$_2$O)$_6$]$^{2+}$ & (R) & low spin complex ion \\
\hline
(IV) & [Fe(H$_2$O)$_6$]$^{2+}$ & (S) & metal ion in 4+ oxidation state \\
\hline
& & (T) & $d^4$ species \\
\hline
\end{tabular}
\end{center}
\begin{multicols}{2}
[Given : Atomic number of Cr = 24, Ru = 44, Fe = 26]

Metal each metal species in LIST-I with their properties in LIST-II, and choose the correct option
\begin{qoptions}
\opt{A}{I$\to$Q,T; II$\to$S,T; III$\to$P,T; IV$\to$Q,R}
\opt{B}{I$\to$R,S; II$\to$P,T; III$\to$P,Q; IV$\to$Q,T}
\opt{C}{I$\to$P,R; II$\to$R,S; III$\to$R,T; IV$\to$P,T}
\opt{D}{I$\to$R,T; II$\to$P,S; III$\to$Q,T; IV$\to$P,Q}
\end{qoptions}

\qn{10} Match the compounds in LIST-I with the observation in LIST-II, and choose the correct option.
\end{multicols}
\begin{center}
\begin{tabular}{|c|c|c|p{6.5cm}|}
\hline
\multicolumn{2}{|c|}{LIST-I} & \multicolumn{2}{c|}{LIST-II} \\
\hline
(I) & Aniline & (P) & Sodium fusion extract of the compound on boiling with FeSO$_4$, followed by acidification with conc. H$_2$SO$_4$, gives Prussian blue color. \\
\hline
(II) & o-Cresol & (Q) & Sodium fusion extract of the compound on treatment with sodium nitroprusside gives blood red color. \\
\hline
(III) & Cysteine & (R) & Addition of the compound to a saturated solution of NaHCO$_3$ results in effervescence. \\
\hline
(IV) & Coprolactam & (S) & The compound reacts with bromine water to give a white precipitate. \\
\hline
& & (T) & Treating the compound with neutral FeCl$_3$ solution produces violet color. \\
\hline
\end{tabular}
\end{center}
\begin{multicols}{2}
\begin{qoptions}
\opt{A}{I$\to$P,Q; II$\to$S; III$\to$Q,R; IV$\to$P}
\opt{B}{I$\to$P; II$\to$R,S; III$\to$R; IV$\to$Q,S}
\opt{C}{I$\to$Q,S; II$\to$P,T; III$\to$P; IV$\to$S}
\opt{D}{I$\to$P,S; II$\to$T; III$\to$Q,R; IV$\to$P}
\end{qoptions}

\sect{Section-II: Numerical Value}

\qn{1} 2 mol of Hg(g) is combusted in a fixed volume bomb calorimeter with excess of O$_2$ at 298 K and 1 atm into HgO(s). During the reaction, temperature increases from 298.0 K to 312.8 K. If heat capacity of the bomb calorimeter and enthalpy of formation of Hg(g) are 20.00 $kJ\,K^{-1}$ and 61.32 $kJ\,mol^{-1}$ at 298 K, respectively, the calculated standard molar enthalpy of formation of HgO(s) at 298 K is X $kJ\,mol^{-1}$. The value of $|X|$ is\underline{\hspace{2em}}.
[Given : Gas constant $R = 8.3\ J\,K^{-1}mol^{-1}$]

\qn{2} The reduction potential ($E^0$, in V) of MnO$_4^-$(aq)/Mn$^{2+}$(aq) is\underline{\hspace{2em}}.
[Given : $E^0_{(MnO_4^-(aq)/MnO_2(s))} = 1.68$ V; $E^0_{(MnO_2(s)/Mn^{2+}(aq))} = 1.21$ V; $E^0_{(Mn^{2+}(aq)/Mn(s))} = -1.03$ V]

\qn{3} A solution is prepared by mixing 0.02 mol each of H$_2$CO$_3$, NaHCO$_3$, Na$_2$CO$_3$, and NaOH in 100 mL of water. pH of the resulting solution is\underline{\hspace{2em}}.
[Given : $pK_{a1}$ and $pK_{a2}$ of H$_2$CO$_3$ are 6.37 and 10.32, respectively ; $\log 2 = 0.30$]

\qn{4} The treatment of an aqueous solution of 7.48 g of Cu(NO$_3$)$_2$ with excess KI results in a brown solution along with the formation of a precipitate. Passing H$_2$S through this brown solution gives another precipitate X. The amount of X (in g) is\underline{\hspace{2em}}.
[Given : Atomic mass of H = 1, N = 14, O = 16, S = 32, K = 39, Cu = 63, I = 127]

\qn{5} Dissolving 2.48 g of white phosphorous in boiling NaOH solution in an inert atmosphere gives a gas $\mathbf{Q}$. The amount of CuSO$_4$ (in g) required to completely consume the gas $\mathbf{Q}$ is\underline{\hspace{2em}}.
[Given : Atomic mass of H = 1, O = 16, Na = 23, P = 31, S = 32, Cu = 63]

\qn{6} Consider the following reaction.
\qfig{p1_chem_s3_q06_1}
On estimation of bromine in 1.00 g of $\mathbf{R}$ using Carius method, the amount of AgBr formed (in g) is\underline{\hspace{2em}}.
[Given : Atomic mass of H = 1, C = 12, O = 16, P = 31, Br = 80, Ag = 108]

\qn{7} The weight percentage of hydrogen in $\mathbf{Q}$, formed in the following reaction sequence, is\underline{\hspace{2em}}.
\qfig{p1_chem_s3_q07_1}
[Given : Atomic mass of H = 1, C = 12, N = 14, O = 16, S = 32, Cl = 35]

\qfullbegin
\qn{8} If the reaction sequence given below is carried out with 15 moles of acetylene, the amount of the product $\mathbf{D}$ formed (in g) is\underline{\hspace{2em}}.
\qfig{p1_chem_s3_q08_2}
\qfullsplit
The yields of $\mathbf{A}$, $\mathbf{B}$, $\mathbf{C}$ and $\mathbf{D}$ are given in parentheses.
[Given : Atomic mass of H = 1, C = 12, O = 16, Cl = 35]
\qfullend

\partsec{3}{MATHEMATICS}
\sect{Section-I (i): One or More Correct}

\qn{1} Consider the equation
\[\int_1^e \dfrac{(\log_e x)^{1/2}}{x\left(k-(\log_e x)^{3/2}\right)^2}dx = 1,\quad k\in(-\infty,0)\cup(1,\infty).\]
Which of the following statements is/are TRUE ?
\begin{qoptions}
\opt{A}{\textbf{No} $k$ satisfies the above equation}
\opt{B}{An integer $k$ satisfies the above equation}
\opt{C}{An irrational number $k$ satisfies the above equation}
\opt{D}{More than one $k$ satisfy the above equation}
\end{qoptions}

\qn{2} Let $a_1,a_2,a_3,\ldots$ be an arithmetic progression with $a_1=7$ and common difference 8. Let $T_1,T_2,T_3,\ldots$ be such that $T_1=5$ and $T_{n+1}-T_n=a_n$ for $n\ge1$. Then, which of the following is/are TRUE ?
\begin{qoptions}
\opt{A}{$T_{10} = 356$}
\opt{B}{$\displaystyle\sum_{k=1}^{20} T_k = 10550$}
\opt{C}{$T_{40} = 6206$}
\opt{D}{$\displaystyle\sum_{k=1}^{30} T_k = 35615$}
\end{qoptions}

\qn{3} Let $P_1$ and $P_2$ be two planes given by
\[P_1: 10x+15y+12z-60=0,\]
\[P_2: -2x+5y+4z-20=0.\]
Which of the following straight lines can be an edge of some tetrahedron whose two faces lie on $P_1$ and $P_2$?
\begin{qoptions}
\opt{A}{$\dfrac{x}{-2}=\dfrac{y-4}{5}=\dfrac{z}{4}$}
\opt{B}{$\dfrac{x-1}{0}=\dfrac{y-1}{0}=\dfrac{z-1}{5}$}
\opt{C}{$\dfrac{x}{1}=\dfrac{y-4}{-2}=\dfrac{z}{3}$}
\opt{D}{$\dfrac{x-6}{-5}=\dfrac{y}{2}=\dfrac{z}{3}$}
\end{qoptions}

\qn{4} Let $S$ be the reflection of a point $Q$ with respect to the plane given by
\[\vec r = -(t+p)\hat\imath + t\hat\jmath + (1+p)\hat k\]
where $t,p$ are real parameters and $\hat\imath,\hat\jmath,\hat k$ are the unit vectors along the three positive coordinate axes. If the position vectors of $Q$ and $S$ are $5\hat\imath+10\hat\jmath+15\hat k$ and $\alpha\hat\imath+\beta\hat\jmath+\gamma\hat k$ respectively, then which of the following is/are TRUE ?
\begin{qoptions}
\opt{A}{$3(\alpha+\beta) = -71$}
\opt{B}{$6(\beta+\gamma) = -71$}
\opt{C}{$3(\gamma+\alpha) = -56$}
\opt{D}{$6(\alpha+\beta+\gamma) = -168$}
\end{qoptions}

\qn{5} Consider the parabola $y^2=4x$. Let $S$ be the focus of the parabola. A pair of tangents drawn to the parabola from the point $P=(-2,1)$ meet the parabola at $P_1$ and $P_2$. Let $Q_1$ and $Q_2$ be points on the lines $SP_1$ and $SP_2$ respectively such that $PQ_1$ is perpendicular to $SP_1$ and $PQ_2$ is perpendicular to $SP_2$. Then, which of the following is/are TRUE ?
\begin{qoptions}
\opt{A}{$SQ_1+SQ_2 = 4$}
\opt{B}{$Q_1Q_2 = \dfrac{2\sqrt{10}}{5}$}
\opt{C}{$PQ_1 = 3$}
\opt{D}{$PQ_2+SQ_1 = 4$}
\end{qoptions}

\qn{6} Let $|M|$ denote the determinant of a square matrix $M$. Let g$:\left[0,\dfrac{\pi}{2}\right]\to\mathbb R$ be the function defined by
\[g(\theta) = \sqrt{f(\theta)-1} + \sqrt{f\left(\dfrac{\pi}{2}-\theta\right)-1}\ \text{where}\]
\par\noindent\resizebox{\linewidth}{!}{$\displaystyle f(\theta) = \dfrac12\begin{vmatrix}1 & \sin\theta & 1\\ -\sin\theta & 1 & \sin\theta\\ -1 & -\sin\theta & 1\end{vmatrix} + \begin{vmatrix}\sin\pi & \cos\left(\theta+\dfrac{\pi}{4}\right) & \tan\left(\theta-\dfrac{\pi}{4}\right)\\ \sin\left(\theta-\dfrac{\pi}{4}\right) & -\cos\dfrac{\pi}{2} & \log_e\left(\dfrac{4}{\pi}\right)\\ \cot\left(\theta+\dfrac{\pi}{4}\right) & \log_e\left(\dfrac{\pi}{4}\right) & \tan\pi\end{vmatrix}$}\par
Let p($x$) be a quadratic polynomial whose roots are the maximum and minimum values of the function g($\theta$), and p(2) $=4-2\sqrt2$. Then, which of the following is/are TRUE ?
\begin{qoptions}
\opt{A}{p$\left(\dfrac{3+\sqrt2}{4}\right)>0$}
\opt{B}{p$\left(\dfrac{1+3\sqrt2}{4}\right)<0$}
\opt{C}{p$\left(\dfrac{5\sqrt2-1}{4}\right)<0$}
\opt{D}{p$\left(\dfrac{5-\sqrt2}{4}\right)>0$}
\end{qoptions}

\sect{Section-I (ii): Matching List (Single Correct)}

\qn{7} Consider the following lists :
\end{multicols}
\begin{center}
\begin{tabular}{|c|p{8cm}|c|l|}
\hline
\multicolumn{2}{|c|}{List-I} & \multicolumn{2}{c|}{List-II} \\
\hline
(I) & $\left\{x\in\left(-\dfrac{2\pi}{3},\dfrac{2\pi}{3}\right)-\{0\}:\cos x+\sin x=1\right\}$ & (P) & has one element \\
\hline
(II) & $\left\{x\in\left[-\dfrac{5\pi}{18},\dfrac{5\pi}{18}\right]:\sqrt3\tan3x=1\right\}$ & (Q) & has two elements \\
\hline
(III) & $\left\{x\in\left(-\dfrac{7\pi}{4},\dfrac{7\pi}{4}\right)-\{0\}:\sin x-\cos x=1\right\}$ & (R) & has three elements \\
\hline
(IV) & $\left\{x\in\left[-\dfrac{6\pi}{5},\dfrac{6\pi}{5}\right]:2\cos(2x)=\sqrt3\right\}$ & (S) & has four elements \\
\hline
& & (T) & has six elements \\
\hline
\end{tabular}
\end{center}
\begin{multicols}{2}
The correct option is :
\begin{qoptions}
\opt{A}{I$\to$P; II$\to$Q; III$\to$S; IV$\to$T}
\opt{B}{I$\to$P; II$\to$P; III$\to$T; IV$\to$R}
\opt{C}{I$\to$Q; II$\to$P; III$\to$T; IV$\to$S}
\opt{D}{I$\to$Q; II$\to$S; III$\to$P; IV$\to$R}
\end{qoptions}

\qn{8} Two players, $P_1$ and $P_2$, play a game against each other. In every round of the game, each player rolls a fair die once, where the six faces of the die have six distinct numbers. Let $x$ and $y$ denote the readings on the die rolled by $P_1$ and $P_2$, respectively. If $x>y$, then $P_1$ scores 5 points and $P_2$ scores 0 point. If $x=y$, then each player scores 2 points. If $x<y$, then $P_1$ scores 0 point and $P_2$ scores 5 points. Let $X_i$ and $Y_i$ be the total scores of $P_1$ and $P_2$, respectively, after playing the $i^{th}$ round.
\end{multicols}
\begin{center}
\begin{tabular}{|c|l|c|l|}
\hline
\multicolumn{2}{|c|}{List-I} & \multicolumn{2}{c|}{List-II} \\
\hline
(I) & Probability of $(X_2\ge Y_2)$ is & (P) & $\dfrac38$ \\
\hline
(II) & Probability of $(X_2>Y_2)$ is & (Q) & $\dfrac{11}{16}$ \\
\hline
(III) & Probability of $(X_3=Y_3)$ is & (R) & $\dfrac{5}{16}$ \\
\hline
(IV) & Probability of $(X_3>Y_3)$ is & (S) & $\dfrac{355}{864}$ \\
\hline
& & (T) & $\dfrac{77}{432}$ \\
\hline
\end{tabular}
\end{center}
\begin{multicols}{2}
The correct option is:
\begin{qoptions}
\opt{A}{I$\to$P; II$\to$R; III$\to$Q; IV$\to$T}
\opt{B}{I$\to$Q; II$\to$R; III$\to$T; IV$\to$T}
\opt{C}{I$\to$P; II$\to$R; III$\to$Q; IV$\to$S}
\opt{D}{I$\to$Q; II$\to$R; III$\to$T; IV$\to$S}
\end{qoptions}

\qn{9} Let $p,q,r$ be nonzero real numbers that are, respectively, the $10^{th}$, $100^{th}$ and $1000^{th}$ terms of a harmonic progression. Consider the system of linear equations
\[x+y+z=1\]
\[10x+100y+1000z=0\]
\[qr\,x+pr\,y+pq\,z=0.\]
\end{multicols}
\begin{center}
\begin{tabular}{|c|p{7cm}|c|l|}
\hline
\multicolumn{2}{|c|}{List-I} & \multicolumn{2}{c|}{List-II} \\
\hline
(I) & If $\dfrac{p}{q}=10$, then the system of linear equations has & (P) & No solution \\
\hline
(II) & If $\dfrac{p}{r}\ne100$, then the system of linear equations has & (Q) & $x=0,\ y=\dfrac{10}{9},\ z=-\dfrac19$ as a solution \\
\hline
(III) & If $\dfrac{q}{r}=10$, then the system of linear equations has & (R) & atleast one solution \\
\hline
(IV) & If $\dfrac{p}{q}\ne10$, then the system of linear equations has & (S) & $x=\dfrac{10}{9},\ y=-\dfrac19,\ z=0$ as a solution \\
\hline
& & (T) & infinitely many solution \\
\hline
\end{tabular}
\end{center}
\begin{multicols}{2}
The correct option is :
\begin{qoptions}
\opt{A}{I$\to$T; II$\to$P; III$\to$S; IV$\to$P}
\opt{B}{I$\to$Q; II$\to$S; III$\to$S; IV$\to$R}
\opt{C}{I$\to$Q; II$\to$R; III$\to$P; IV$\to$R}
\opt{D}{I$\to$T; II$\to$S; III$\to$P; IV$\to$P}
\end{qoptions}

\qn{10} Consider the ellipse
\[\dfrac{x^2}{4}+\dfrac{y^2}{3}=1.\]
Let H($\alpha$, 0), $0<\alpha<2$, be a point. A straight line drawn through $H$ parallel to the $y$-axis crosses the ellipse and its auxiliary circle at points $E$ and $F$ respectively, in the first quadrant. The tangent to the ellipse at the point $E$ intersects the positive $x$-axis at a point $G$. Suppose the straight line joining $F$ and the origin makes an angle $\theta$ with the positive $x$-axis.
\end{multicols}
\begin{center}
\begin{tabular}{|c|p{7cm}|c|l|}
\hline
\multicolumn{2}{|c|}{List-I} & \multicolumn{2}{c|}{List-II} \\
\hline
(I) & If $\theta=\dfrac{\pi}{12}$, then the area of the triangle $FGH$ is & (P) & $\dfrac{1}{2\sqrt3}$ \\
\hline
(II) & If $\theta=\dfrac{\pi}{4}$, then the area of the triangle $FGH$ is & (Q) & $\dfrac{(\sqrt3-1)^4}{4}$ \\
\hline
(III) & If $\theta=\dfrac{\pi}{3}$, then the area of the triangle $FGH$ is & (R) & $\dfrac{(2-\sqrt3)^2}{2}$ \\
\hline
(IV) & If $\theta=\dfrac{\pi}{6}$, then the area of the triangle $FGH$ is & (S) & $\dfrac{3\sqrt3}{2}$ \\
\hline
& & (T) & 1 \\
\hline
\end{tabular}
\end{center}
\begin{multicols}{2}
The correct option is :
\begin{qoptions}
\opt{A}{I$\to$R; II$\to$S; III$\to$Q; IV$\to$S}
\opt{B}{I$\to$R; II$\to$T; III$\to$S; IV$\to$P}
\opt{C}{I$\to$Q; II$\to$T; III$\to$S; IV$\to$P}
\opt{D}{I$\to$Q; II$\to$S; III$\to$Q; IV$\to$S}
\end{qoptions}

\sect{Section-II: Numerical Value}

\qn{1} Considering only the principal values of the inverse trigonometric functions, the value of
\[\dfrac32\cos^{-1}\sqrt{\dfrac{2}{2+\pi^2}} + \dfrac14\sin^{-1}\dfrac{2\sqrt2\,\pi}{2+\pi^2} + \tan^{-1}\dfrac{\sqrt2}{\pi}\]
is $\dfrac{\lambda\pi}{3}$ then value of $\lambda$ is\underline{\hspace{2em}}.

\qn{2} Let $\alpha$ be a positive real number. Let $f:\mathbb R\to\mathbb R$ and g$:(\alpha,\infty)\to\mathbb R$ be the functions defined by
\[f(x) = \sin\left(\dfrac{\pi x}{2}\right)\ \text{and}\ g(x) = \dfrac{\log_e\left(\sqrt x-\sqrt\alpha\right)}{\log_e\left(e^{\sqrt x}-e^{\sqrt\alpha}\right)}.\]
Then the value of $\displaystyle\lim_{x\to\alpha^+} f(g(x))$ is\underline{\hspace{2em}}.

\qn{3} In a study about a pandemic, data of 900 persons was collected. It was found that

190 persons had symptom of fever,

220 persons had symptom of cough,

220 persons had symptom of breathing problem,

330 persons had symptom of fever or cough or both,

350 persons had symptom of cough or breathing problem or both,

340 persons had symptom of fever or breathing problem or both,

30 persons had all three symptoms (fever, cough and breathing problem).

If a person is chosen randomly from these 900 persons, then the probability that the person is suffering from exactly one symptom is\underline{\hspace{2em}}.

\qn{4} Let $z$ be a complex number with non-zero imaginary part. If
\[\dfrac{1+3z+4z^2}{1-3z+4z^2}\]
is a real number, then the value of $|z|^2$ is\underline{\hspace{2em}}.

\qn{5} Let $\bar z$ denote the complex conjugate of a complex number $z$ and let $i=\sqrt{-1}$. In the set of complex numbers, the number of distinct roots of the equation $\bar z-z^2=i(\bar z+z^2)$ is\underline{\hspace{2em}} (where Re(z) $\ne0$)

\qn{6} Let $l_1,l_2,\ldots,l_{100}$ be consecutive terms of an arithmetic progression with common difference $d_1$, and let $w_1,w_2,\ldots,w_{100}$ be consecutive terms of another arithmetic progression with common difference $d_2$, where $d_1d_2=1$. For each $i=1,2,\ldots,100$, let $R_i$ be a rectangle with length $l_i$, width $w_i$ and area $A_i$. If $A_{51}-A_{50}=100$, then the value of $\dfrac{A_{100}-A_{90}}{5}$ is\underline{\hspace{2em}}.

\qn{7} If the number of 4-digit integers in the closed interval [2032, 4482] formed by using the digits 0, 2, 3, 4, 6, 7 is $\mu$ then value of $\left(\dfrac{\mu-500}{2}\right)$ is\underline{\hspace{2em}}.

\qn{8} Let $ABC$ be the triangle with $AB=1$, $AC=3$ and $\angle BAC=\dfrac{\pi}{2}$. If a circle of radius $k>0$ touches the sides $AB$, $AC$ and also touches internally the circumcircle of the triangle $ABC$, then the value of [k] is\underline{\hspace{2em}} (where [.] denotes greatest integer function)

\papersec{2}
\partsec{1}{PHYSICS}
\sect{Section-I (i): Single Correct}

\qn{1} A flat surface of a thin uniform disk $A$ of radius $R$ is glued to a horizontal table. Another thin uniform disk $B$ of mass $M$ and with the radius $2R$ rolls without slipping on the circumference of $A$, as shown in the figure. A flat surface of $B$ also lies on the plane of the table. The center of mass of $B$ has fixed angular speed $\omega$ about the vertical axis passing through the center of $A$. The angular momentum of $B$ is $nM\omega R^2$ with respect to the center of $A$. Which of the following is the value of $n$?
\qfig{p2_phy_s1_q01_1}
\begin{qoptions}
\opt{A}{15}
\opt{B}{5}
\opt{C}{$\dfrac72$}
\opt{D}{$\dfrac92$}
\end{qoptions}

\qn{2} When light of a given wavelength is incident on a metallic surface, the minimum potential needed to stop the emitted photoelectrons is 6.0 $V$. This potential drops to 0.6 $V$ if another source with wavelength two times that of the first one and intensity double of the first one is used. What are the wavelength of the first source and the work function of the metal, respectively? [Take $=\dfrac{hc}{e}=1.24\times10^{-6}\,Jm\,C^{-1}$.]
\begin{qoptions}
\opt{A}{$2.29\times10^{-7}$ m, 4.80 eV}
\opt{B}{$1.15\times10^{-7}$ m, 4.80 eV}
\opt{C}{$3.78\times10^{-7}$ m, 5.60 eV}
\opt{D}{$3.78\times10^{-7}$ m, 1.20 eV}
\end{qoptions}

\qn{3} Area of the cross-section of a wire is measured using a screw gauge. The pitch of the main scale is 0.5 mm. The circular scale has 100 divisions and for one full rotation of the circular scale, the main scale shifts by two divisions. The measured readings are listed below.
\end{multicols}
\begin{center}
\begin{tabular}{|p{4.5cm}|p{4cm}|p{4cm}|}
\hline
Measurement condition & Main scale reading & Circular scale reading \\
\hline
Two arms of gauge touching each other without wire & 0 division & 4 division \\
\hline
Attempt-1: With wire & 4 divisions & 20 divisions \\
\hline
Attempt-2: With wire & 4 divisions & 16 divisions \\
\hline
\end{tabular}
\end{center}
\begin{multicols}{2}
What are the diameter and cross-sectional area of the wire measured using the screw gauge?
\begin{qoptions}
\opt{A}{$1.07\pm0.02$ mm, $\pi(1.14\pm0.02)$ mm$^2$}
\opt{B}{$1.07\pm0.01$ mm, $\pi(0.29\pm0.005)$ mm$^2$}
\opt{C}{$2.07\pm0.02$ mm, $\pi(1.07\pm0.02)$ mm$^2$}
\opt{D}{$2.14\pm0.02$ mm, $\pi(1.07\pm0.02)$ mm$^2$}
\end{qoptions}

\qn{4} Which one of the following options represents the magnetic field $\vec B$ at O due to the current flowing in the given wire segments lying on the $xy$ plane?
\qfig{p2_phy_s1_q04_1}
\begin{qoptions}
\opt{A}{$\vec B = \dfrac{-\mu_0 I}{L}\left(\dfrac32+\dfrac{1}{4\sqrt2\pi}\right)\hat k$}
\opt{B}{$\vec B = -\dfrac{\mu_0 I}{L}\left(\dfrac32+\dfrac{1}{2\sqrt2\pi}\right)\hat k$}
\opt{C}{$\vec B = \dfrac{-\mu_0 I}{L}\left(1+\dfrac{1}{4\sqrt2\pi}\right)\hat k$}
\opt{D}{$\vec B = \dfrac{-\mu_0 I}{L}\left(1+\dfrac{1}{4\pi}\right)\hat k$}
\end{qoptions}

\sect{Section-I (ii): One or More Correct}

\qn{5} In the figure, the inner (shaded) region $A$ represents a sphere of radius $r_A=1$, within which the electrostatic charge density varies with the radial distance $r$ from the center as $\rho_A=kr$, where $k$ is positive. In the spherical shell $B$ of outer radius $r_B$, the electrostatic charge density varies as $\rho_B=\dfrac{k}{r}$. Assume that dimensions are taken care of. All physical quantities are in their SI units.
\qfig{p2_phy_s2_q05_1}
Which of the following statement(s) is(are) correct?
\begin{qoptions}
\opt{A}{It is not possible that electric field is zero everywhere outside.}
\opt{B}{If $r_B=\dfrac32$, then the electric potential just outside $B$ is $\dfrac{7k}{12\epsilon_0}$.}
\opt{C}{If $r_B=2$, then the total charge of the configuration is $7\pi k$.}
\opt{D}{If $r_B=\dfrac52$, then the magnitude of the electric field just outside $B$ is $\dfrac{13\pi k}{\epsilon_0}$.}
\end{qoptions}

\qn{6} In Circuit-1 and Circuit-2 shown in the figures, $R_1=1\,\Omega$, $R_2=2\,\Omega$ and $R_3=3\,\Omega$. $P_1$ and $P_2$ are the power dissipations in Circuit-1 and Circuit-2 when the switches $S_1$ and $S_2$ are in open conditions, respectively. $Q_1$ and $Q_2$ are the power dissipations in Circuit-1 and Circuit-2 when the switches $S_1$ and $S_2$ are in closed conditions, respectively.
\qfig{p2_phy_s2_q06_2}
\qfig{p2_phy_s2_q06_1}
Which of the following statement(s) is(are) correct?
\begin{qoptions}
\opt{A}{When a voltage source of 12 $V$ is connected across A and B in both circuits, $P_1<P_2$.}
\opt{B}{When a constant current source of 3 $Amp$ is connected across A and B in both circuits, $P_1>P_2$.}
\opt{C}{When a voltage source of 4 $V$ is connected across A and B in Circuit-1, $Q_1>P_1$.}
\opt{D}{When a constant current source of 3 $Amp$ is connected across A and B in both circuits, $Q_2<Q_1$}
\end{qoptions}

\qn{7} A bubble has surface tension $S$. The ideal gas inside the bubble has ratio of specific heats $\gamma=\dfrac53$. The bubble is exposed to the atmosphere and it always retains its spherical shape. When the atmospheric pressure is P$_{a1}$, the radius of the bubble is found to be $r_1$ and the temperature of the enclosed gas is $T_1$. When the atmospheric pressure is P$_{a2}$, the radius of the bubble and the temperature of the enclosed gas are $r_2$ and $T_2$, respectively. Which of the following statement(s) is(are) correct?
\begin{qoptions}
\opt{A}{If the surface of the bubble is a perfect heat insulator, then $\left(\dfrac{r_1}{r_2}\right)^5 = \dfrac{P_{a2}+\dfrac{2S}{r_2}}{P_{a1}+\dfrac{2S}{r_1}}$}
\opt{B}{If the surface of the bubble is a perfect heat insulator, then the total internal energy of the bubble including its surface energy does not change with the external atmospheric pressure.}
\opt{C}{If the surface of the bubble is a perfect heat conductor and the change in atmospheric temperature is negligible, then $\left(\dfrac{r_1}{r_2}\right)^3 = \dfrac{P_{a2}+\dfrac{4S}{r_2}}{P_{a1}+\dfrac{4S}{r_1}}$.}
\opt{D}{If the surface of the bubble is a perfect heat insulator, then $\left(\dfrac{T_2}{T_1}\right)^{5/2} = \dfrac{P_{a2}+\dfrac{4S}{r_2}}{P_{a1}+\dfrac{4S}{r_1}}$.}
\end{qoptions}

\qn{8} A disk of radius $R$ with uniform positive charge density $\sigma$ is placed on the $xy$ plane with its center at the origin. The Coulomb potential along the $z$-axis is V$(z)=\dfrac{\sigma}{2\epsilon_0}\left(\sqrt{R^2+z^2}-z\right)$. A particle of positive charge $q$ is placed initially at rest at a point on the $z$ axis with $z=z_0$ and $z_0>0$. In addition to the Coulomb force, the particle experiences a vertical force $\vec F=-c\hat k$ with $c>0$. Let $\beta=\dfrac{2c\epsilon_0}{q\sigma}$. Which of the following statement(s) is(are) correct? (Ignore gravity)
\begin{qoptions}
\opt{A}{For $\beta=\dfrac14$ and $z_0=\dfrac{25}{7}R$, the particle reaches the origin.}
\opt{B}{For $\beta=\dfrac14$ and $z_0=\dfrac37R$, the particle reaches the origin.}
\opt{C}{For $\beta=\dfrac14$ and $z_0=\dfrac{R}{\sqrt3}$, the particle returns back to $z=z_0$.}
\opt{D}{For $\beta>1$ and $z_0>0$, the particle always reaches the origin.}
\end{qoptions}

\qn{9} A double slit setup is shown in the figure. One of the slits is in medium 2 of refractive index $n_2$. The other slit is at the interface of this medium with another medium 1 of refractive index $n_1(\ne n_2)$. The line joining the slits is perpendicular to the interface and the distance between the slits is $d$. The slit widths are much smaller than $d$. A monochromatic parallel beam of light is incident on the slits from medium 1. A detector is placed in medium 2 at a large distance from the slits, and at an angle $\theta$ from the line joining them, so that $\theta$ equals the angle of refraction of the beam. Consider two approximately parallel rays from the slits received by the detector. Which of the following statement(s) is (are) correct?
\qfig{p2_phy_s2_q09_1}
\begin{qoptions}
\opt{A}{The phase difference between the two rays is independent of $d$.}
\opt{B}{The two rays interfere constructively at the detector.}
\opt{C}{The phase difference between the two rays depends on $n_1$ but is independent of $n_2$.}
\opt{D}{The phase difference between the two rays vanishes only for certain values of $d$ and the angle of incidence of the beam, with $\theta$ being the corresponding angle of refraction.}
\end{qoptions}

\qn{10} In the given $P$-$V$ diagram, a monoatomic gas $\left(\gamma=\dfrac53\right)$ is first compressed adiabatically from state $A$ to state $B$. Then it expands isothermally from state $B$ to state $C$. [Given: $\left(\dfrac13\right)^{0.6}\simeq0.5$, $\ln2\simeq0.7$].
\qfig{p2_phy_s2_q10_1}
Which of the following statement(s) is(are) correct?
\begin{qoptions}
\opt{A}{The modulus of the total work done in the process $A\to B\to C$ is 144 $kJ$.}
\opt{B}{The modulus of the work done in the process $B\to C$ is 84 $kJ$.}
\opt{C}{The modulus of the work done in the process $A\to B$ is 60 $kJ$.}
\opt{D}{The modulus of the work done in the process $C\to A$ is zero.}
\end{qoptions}

\sect{Section-III: Single Digit Integer (0-9)}

\qn{1} A particle of mass 1 $kg$ is subjected to a force which depends on the position as $\vec F=-k(x\hat\imath+y\hat\jmath)$ $kg\,ms^{-2}$ with $k=1\ kg\,s^{-2}$. At time $t=0$, the particle's position is $\vec r=\left(\dfrac{1}{\sqrt2}\hat\imath+\sqrt2\hat\jmath\right)$ m and its velocity $\vec v=\left(-\sqrt2\hat\imath+\sqrt2\hat\jmath+\dfrac{4}{\pi}\hat k\right)\,ms^{-1}$. Let $v_x$ and $v_y$ denote the $x$ and $y$ components of the particle's velocity, respectively. \textbf{Ignore gravity.} When $z=0.5$ m, the value of $(xv_y-yv_x)$ is\underline{\hspace{2em}} $m^2s^{-1}$.

\qn{2} In a radioactive decay chain reaction, ${}^{230}_{90}$Th nucleus decays into ${}^{214}_{84}$Po nucleus. The sum of the number of $\alpha$ and number of $\beta^-$ particles emitted in this process is\underline{\hspace{2em}}.

\qn{3} Two resistances $R_1=X\,\Omega$ and $R_2=0.4\,\Omega$ are connected to a wire AB of uniform resistivity, as shown in the figure. The radius of the wire varies linearly along its axis from 0.2 mm at A to 1 mm at B. A galvanometer (G) connected to the center of the wire, 50 cm from each end along its axis, shows zero deflection when A and B are connected to a battery. The value of X (in $\Omega$) is\underline{\hspace{2em}}.
\qfig{p2_phy_s3_q03_1}

\qn{4} In a particular system of units, a physical quantity can be expressed in terms of the electric charge e, electron mass m$_e$, Planck's constant h, and Coulomb's constant $k=\dfrac{1}{4\pi\epsilon_0}$, where $\epsilon_0$ is the permittivity of vacuum. In terms of these physical constants, the dimension of the magnetic field is $[B]=[e]^\alpha[m_e]^\beta[h]^\gamma[k]^\delta$. The value of $\alpha+\beta+\gamma+\delta$ is\underline{\hspace{2em}}.

\qn{5} Consider a configuration of n identical units, each consisting of three layers. The first layer is a column of air of height h $=\dfrac13$ cm, and the second and third layers are of equal thickness d $=\dfrac{\sqrt3-1}{2}$ cm, and refractive indices $\mu_1=\sqrt{\dfrac32}$ and $\mu_2=\sqrt3$, respectively. A light source O is placed on the top of the first unit, as shown in the figure. A ray of light from O is incident on the second layer of the first unit at an angle of $\theta=60^\circ$ to the normal. For a specific value of n, the ray of light emerges from the bottom of the configuration at a distance $\ell=\dfrac{8}{\sqrt3}$ cm, as shown in the figure. The value of n is\underline{\hspace{2em}}.
\qfig{p2_phy_s3_q05_1}

\qn{6} A charge $q$ is surrounded by a closed surface consisting of an inverted cone of height $h$ and base radius $R$, and a half ellipsoid of base radius $R$ as shown in the figure. The electric flux through the conical surface is $\dfrac{nq}{6\epsilon_0}$ (in SI units). The value of $n$ is\underline{\hspace{2em}}.
\qfig{p2_phy_s3_q06_1}

\qn{7} On a frictionless horizontal plane, a bob of mass m $=0.1$ kg is attached to a spring with natural length $\ell_0=0.1$ m. The spring constant is $k_1=0.001\,Nm^{-1}$ when the length of the spring $\ell>\ell_0$ and is $k_2=0.004\,Nm^{-1}$ when $\ell<\ell_0$. Initially the bob is released from $\ell=0.15$ m. Assume that Hooke's law remains valid throughout the motion. If the time period of the full oscillation is T $=(3n\pi)$ s, then the integer closest to n is\underline{\hspace{2em}}.

\qn{8} An object and a concave mirror of focal length f $=10$ cm both move along the principal axis of the mirror with constant speeds. The object moves with speed V$_0=6\,cm\,s^{-1}$ towards the mirror with respect to a laboratory frame. The distance between the object and the mirror at a given moment is denoted by u. When u $=20$ cm, the speed of the mirror V$_m$ is such that the image is instantaneously at rest with respect to the laboratory frame, and the object forms a real image. The magnitude of V$_m$ is\underline{\hspace{2em}} $cm\,s^{-1}$.
\qfig{p2_phy_s3_q08_1}

\partsec{2}{CHEMISTRY}
\sect{Section-I (i): Single Correct}

\qn{1} In a crystalline solid atoms 'X' occupy the fcc lattice sites as well as alternate tetrahedral voids in the same lattice. Calculate \% vacant space in the resulting solid ?
\begin{qoptions}
\opt{A}{74}
\opt{B}{34}
\opt{C}{26}
\opt{D}{66}
\end{qoptions}

\qn{2} The reaction of HClO$_3$ with HCl gives a paramagnetic gas, which upon reaction with O$_3$ produces
\begin{qoptions}
\opt{A}{Cl$_2$O}
\opt{B}{ClO$_2$}
\opt{C}{Cl$_2$O$_7$}
\opt{D}{Cl$_2$O$_6$}
\end{qoptions}

\qn{3} The reaction Pb(NO$_3$)$_2$ and NaCl in water produces a precipitate that dissolves upon the addition of HCl of appropriate concentration. The dissolution of the precipitate is due to the formation of
\begin{qoptions}
\opt{A}{PbCl$_2$}
\opt{B}{PbCl$_4$}
\opt{C}{[PbCl$_6$]$^{2-}$}
\opt{D}{[PbCl$_4$]$^{2-}$}
\end{qoptions}

\qn{4} Treatment of D-glucose with aqueous NaOH results in a mixture of monosaccharides, which are
\begin{qoptions}
\opt{A}{\qfigi{p2_chem_s1_q04_A}}
\opt{B}{\qfigi{p2_chem_s1_q04_B}}
\opt{C}{\qfigi{p2_chem_s1_q04_C}}
\opt{D}{\qfigi{p2_chem_s1_q04_D}}
\end{qoptions}

\sect{Section-I (ii): One or More Correct}

\qn{5} To check the principle of multiple proportions, a series of pure binary compounds (P$_m$Q$_n$) were analyzed and their composition is tabulated below. The correct option(s) is(are)
\begin{center}
\begin{tabular}{|c|c|c|}
\hline
Compound & Weight \% of P & Weight \% of Q \\
\hline
1 & 50 & 50 \\
\hline
2 & 44.4 & 55.6 \\
\hline
3 & 40 & 60 \\
\hline
\end{tabular}
\end{center}
\begin{qoptions}
\opt{A}{If empirical formula of compound 3 is P$_3$Q$_4$, then the empirical formula of compound 2 is P$_3$Q$_5$.}
\opt{B}{If empirical formula of compound 3 is P$_3$Q$_2$ and atomic weight of element P is 20, then the atomic weight of Q is 45.}
\opt{C}{If empirical formula of compound 2 is PQ, then the empirical formula of the compound 1 is P$_5$Q$_4$.}
\opt{D}{If atomic weight of P and Q are 70 and 35, respectively, then the empirical formula of compound 1 is PQ$_2$.}
\end{qoptions}

\qn{6} The correct option(s) about entropy (S) is(are)
[R = gas constant, F = Faraday constant, T = Temperature]
\begin{qoptions}
\opt{A}{For the reaction, M(s) $+ 2H^+(aq)\longrightarrow H_2(g)+M^{2+}(aq)$, if $\dfrac{dE_{cell}}{dT}=\dfrac{R}{F}$, then the entropy change of the reaction is 2R (assume that entropy and internal energy changes are temperature independent).}
\opt{B}{The cell reaction, Pt(s) $\mid$ H$_2(g,$1bar$)\mid$ H$^+(aq,$0.01M$)\parallel$ H$^+(aq,$0.1M$)\mid$ H$_2(g,$1bar$)\mid$ Pt(s), is an enthalpy driven process.}
\opt{C}{For racemization of an optically active compound, $\Delta S<0$.}
\opt{D}{$\Delta S>0$, for [Ni(H$_2$O)$_6$]$^{2+}+3$ en $\longrightarrow$ [Ni(en)$_3$]$^{2+}+6$H$_2$O (where en = ethylenediamine).}
\end{qoptions}

\qn{7} The species which react(s) with NH$_3$ to give boron nitride (BN) is(are)
\begin{qoptions}
\opt{A}{HBF$_4$}
\opt{B}{B$_2$H$_6$}
\opt{C}{B$_2$O$_3$}
\opt{D}{B}
\end{qoptions}

\qn{8} The correct option(s) related to the extraction of iron from its ore in the blast furnace operating in the temperature range $900-1500$ K is(are)
\begin{qoptions}
\opt{A}{Limestone is used to remove silicate impurity.}
\opt{B}{Pig iron obtained from blast furnace contains about 4\% carbon.}
\opt{C}{Coke (C) converts CO to CO$_2$.}
\opt{D}{Exhaust gases consist of NO$_2$ and CO.}
\end{qoptions}

\qn{9} Considering the following reaction sequence, the correct statement(s) is(are)
\qfig{p2_chem_s2_q09_1}
\begin{qoptions}
\opt{A}{Compound $\mathbf{R}$ involved in Rosenmund reaction and give aldehyde.}
\opt{B}{Compound $\mathbf{S}$ show Aldol condensation reaction.}
\opt{C}{Compounds $\mathbf{P}$ and $\mathbf{S}$ react with hydroxylamine to give the corresponding oximes.}
\opt{D}{Compound $\mathbf{R}$ reacts with dialkylcadmium to give the corresponding tertiary alcohol.}
\end{qoptions}

\qn{10} Among the following, the correct statement(s) about polymers is(are)
\begin{qoptions}
\opt{A}{The polymerization of chloroprene gives natural rubber.}
\opt{B}{Teflon is prepared from tetrafluoroethene by heating it with persulphate catalyst at high pressures.}
\opt{C}{Glyptal polymer is condensation polymer}
\opt{D}{Ethene at 350-570 K temperature and 1000-2000 atm pressure in the presence of a peroxide initiator yields high density polythene.}
\end{qoptions}

\sect{Section-III: Single Digit Integer (0-9)}

\qn{1} Concentration of H$_2$SO$_4$ and Na$_2$SO$_4$ in a solution is 1 M and $1.8\times10^{-2}$ M, respectively. Molar solubility of PbSO$_4$ in the same solution is $X\times10^{-Y}$ M (expressed in scientific notation). The value of Y is\underline{\hspace{2em}}.
[Given: Solubility product of PbSO$_4$ ($K_{sp}$) $=1.6\times10^{-8}$. For H$_2$SO$_4$, $K_{a1}$ is very large and $K_{a2}=1.2\times10^{-2}$]

\qn{2} An aqueous solution is prepared by dissolving 0.1 mol of an ionic salt in 1.8 kg of water at $35^\circ$C. The salt remains 90\% dissociated in the solution. The vapour pressure of the solution is 59.832 mm of Hg. Vapor pressure of water at $35^\circ$C is 60.000 mm of Hg. The number of ions present per formula unit of the ionic salt is\underline{\hspace{2em}}.

\qn{3} Consider the strong electrolytes Z$_m$X$_n$, U$_m$Y$_p$ and V$_m$X$_n$. Limiting molar conductivity ($\Lambda^0$) of U$_m$Y$_p$ and V$_m$X$_n$ are 350 and 440 $S\,cm^2\,mol^{-1}$, respectively. The value of (m + n + p) is\underline{\hspace{2em}}.
Given :
\end{multicols}
\begin{center}
\begin{tabular}{|l|c|c|c|c|c|}
\hline
Ion & Z$^{n+}$ & U$^{p+}$ & V$^{n+}$ & X$^{m-}$ & Y$^{m-}$ \\
\hline
$\lambda^0$($S\,cm^2\,mol^{-1}$) & 50.0 & 25.0 & 100.0 & 80.0 & 100.0 \\
\hline
\end{tabular}
\end{center}
\begin{multicols}{2}
$\lambda^0$ is the limiting molar conductivity of ions

The plot of molar conductivity ($\Lambda$) of Z$_m$X$_n$ \textit{vs} $c^{1/2}$ is given below.
\qfig{p2_chem_s3_q03_2}

\qn{4} The reaction of XeF$_4$ and O$_2$F$_2$ gives a Xe compound $\mathbf{P}$. The number of moles of HF produced by the complete hydrolysis of 1 mol of $\mathbf{P}$ is\underline{\hspace{2em}}.

\qn{5} Thermal decomposition of Zn(NO$_3$)$_2$ produces two paramagnetic gases. The total number of bonding molecular orbitals field with electrons of the gas that has the higher number of unpaired electrons is\underline{\hspace{2em}}.

\qn{6} The number of isomeric tetraenes (\textbf{NOT} containing $sp$-hybridized carbon atoms) that can be formed from the following reaction sequence is\underline{\hspace{2em}}.
\qfig{p2_chem_s3_q06_1}

\qn{7} The number of -CH$_2$- (methylene) groups in the product formed from the following reaction sequence is\underline{\hspace{2em}}.
1. O$_3$, Zn/H$_2$O
\qfig{p2_chem_s3_q07_2}

\qn{8} The total number of chiral molecules formed from one molecule of $\mathbf{P}$ on complete ozonolysis (O$_3$, Zn/H$_2$O) is\underline{\hspace{2em}}.
\qfig{p2_chem_s3_q08_1}

\partsec{3}{MATHEMATICS}
\sect{Section-I (i): Single Correct}

\qn{1} Consider 3 boxes, where each box contains 3 red balls and 2 blue balls. Assume that all 15 balls are distinct. In how many different ways can 8 balls be chosen from these 3 boxes so that from each box at least one red ball and at least one blue ball are chosen ?
\begin{qoptions}
\opt{A}{1998}
\opt{B}{2178}
\opt{C}{1209}
\opt{D}{1568}
\end{qoptions}

\qn{2} If $M=\begin{pmatrix}\dfrac52 & \dfrac32\\[4pt] -\dfrac32 & -\dfrac12\end{pmatrix}$, then which of the following matrices is equal to $M^{2024}$?
\begin{qoptions}
\opt{A}{$\begin{pmatrix}3037 & 3036\\ -3036 & -3035\end{pmatrix}$}
\opt{B}{$\begin{pmatrix}3037 & -3036\\ 3036 & -3035\end{pmatrix}$}
\opt{C}{$\begin{pmatrix}3033 & 3032\\ -3032 & -3031\end{pmatrix}$}
\opt{D}{$\begin{pmatrix}3032 & 3031\\ -3031 & -3030\end{pmatrix}$}
\end{qoptions}

\qn{3} Suppose that

Box-I contains 8 red, 3 blue and 5 green balls,

Box-II contains 24 red, 9 blue and 15 green balls,

Box-III contains 1 blue, 12 green and 3 yellow balls,

Box-IV contains 10 green, 16 orange and 6 white balls.

A ball is chosen randomly from Box-I ; call this ball $b$. If $b$ is red then a ball is chosen randomly from Box-II, if $b$ is blue then a ball is chosen randomly from Box-III, and if $b$ is green then a ball chosen randomly from Box-IV. The conditional probability of the event `one of the chosen balls is orange' given that the event `at least one of the chosen balls is green' has happened, is equal to
\begin{qoptions}
\opt{A}{$\dfrac{15}{256}$}
\opt{B}{$\dfrac{10}{39}$}
\opt{C}{$\dfrac{5}{52}$}
\opt{D}{$\dfrac18$}
\end{qoptions}

\qn{4} For positive integer n, define
\par\noindent\resizebox{\linewidth}{!}{$\displaystyle f(n) = n + \dfrac{16+5n-5n^2}{4n+5n^2} + \dfrac{32+n-5n^2}{8n+5n^2} + \dfrac{48-3n-5n^2}{12n+5n^2} + \ldots + \dfrac{25n-5n^2}{9n^2}$}\par
Then, the value of $\displaystyle\lim_{n\to\infty} f(n)$ is equal to
\begin{qoptions}
\opt{A}{$4-\dfrac{11}{4}\log_e\left(\dfrac95\right)$}
\opt{B}{$4-\dfrac34\log_e\left(\dfrac75\right)$}
\opt{C}{$4+\dfrac43\log_e\left(\dfrac73\right)$}
\opt{D}{$3+\dfrac{11}{4}\log\dfrac95$}
\end{qoptions}

\sect{Section-I (ii): One or More Correct}

\qn{5} Let PQRS be a quadrilateral in a plane, where QR $=\sqrt3$, $\angle$PQR $=\angle$QRS $=70^\circ$, $\angle$PQS $=15^\circ$ and $\angle$PRS $=40^\circ$. If $\angle$RPS $=\theta^\circ$, PQ $=\alpha$ and PS $=\beta$, then the interval(s) that does \textbf{NOT} contain(s) the value of $4\alpha\beta\sin\theta^\circ$ is/are ?
\begin{qoptions}
\opt{A}{$(0,\sqrt2)$}
\opt{B}{$(1,2)$}
\opt{C}{$(\sqrt2,3)$}
\opt{D}{$(2\sqrt2,3\sqrt2)$}
\end{qoptions}

\qn{6} Let
\[\alpha = \sum_{k=1}^{\infty}\cos^{2k}\left(\dfrac{\pi}{3}\right)\]
Let $g:[0,1]\to\mathbb R$ be the function defined by
\[g(x) = 3^{\alpha x}+3^{\alpha(1-x)}\]
Then, which of the following statements is/are TRUE?
\begin{qoptions}
\opt{A}{The minimum value of $g(x)$ is $3^{7/6}$}
\opt{B}{The maximum value of $g(x)$ is $1+3^{1/3}$}
\opt{C}{The function $g(x)$ attains its maximum at more than one point}
\opt{D}{The function $g(x)$ attains its minimum at more than one point}
\end{qoptions}

\qn{7} Let $\bar z$ denote the complex conjugate of a complex number $z$. If $z$ is a non-zero complex number for which both real and imaginary parts of
\[(\bar z)^2+\dfrac{1}{z^2}\]
are integers, then which of the following is/are possible value(s) of $|z|$?
\begin{qoptions}
\opt{A}{$\left(\dfrac{3+\sqrt5}{2}\right)^{1/4}$}
\opt{B}{$\left(\dfrac{11+3\sqrt{13}}{2}\right)^{1/4}$}
\opt{C}{$\left(\dfrac{9+\sqrt{65}}{4}\right)^{1/4}$}
\opt{D}{$\left(\dfrac{7+\sqrt{13}}{6}\right)^{1/4}$}
\end{qoptions}

\qn{8} Let $\hat\imath,\hat\jmath$ and $\hat k$ be the unit vectors along the three positive coordinate axes. Let
\[\vec a = \hat\imath+2\hat\jmath-2\hat k,\]
\[\vec b = \hat\imath+b_2\hat\jmath+b_3\hat k,\quad b_2,b_3\in\mathbb R,\]
\[\vec c = c_1\hat\imath+c_2\hat\jmath+c_3\hat k,\quad c_1,c_2,c_3\in\mathbb R\]
be three vectors such that $b_2b_3>0$, $\vec a\cdot\vec b=0$ and
\[\begin{pmatrix}0 & -c_3 & c_2\\ c_3 & 0 & -c_1\\ -c_2 & c_1 & 0\end{pmatrix}\begin{pmatrix}1\\ b_2\\ b_3\end{pmatrix} = \begin{pmatrix}1-c_1\\ 2-c_2\\ -2-c_3\end{pmatrix}.\]
Then, which of the following is/are TRUE ?
\begin{qoptions}
\opt{A}{$\vec a\cdot\vec c = 0$}
\opt{B}{$\vec b\cdot\vec c = 0$}
\opt{C}{$|\vec b|>\dfrac{\sqrt5}{2}$}
\opt{D}{$|\vec c|\le3$}
\end{qoptions}

\qn{9} For $x\in\mathbb R$, let the function $y(x)$ be the solution of the differential equation
\[\dfrac{dy}{dx}+12y = \sin\left(\dfrac{\pi}{12}x\right),\quad y(0)=0.\]
Then, which of the following statements is/are TRUE?
\begin{qoptions}
\opt{A}{$y(x)$ is an increasing function}
\opt{B}{$y(x)$ is a decreasing function}
\opt{C}{There exists a real number $\beta$ such that the line $y=\beta$ intersects the curve $y=y(x)$ at infinitely many points}
\opt{D}{$y(x)$ is a periodic function}
\end{qoptions}

\qn{10} Let G be a circle of radius $R>0$. Let $G_1,G_2,\ldots,G_n$ be n circles of equal radius $r>0$. Suppose each of the n circles $G_1,G_2,\ldots,G_n$ touches the circle G externally. Also, for $i=1,2,\ldots,n-1$, the circle $G_i$ touches $G_{i+1}$ externally, and $G_n$ touches $G_1$ externally. Then, which of the following statements is/are TRUE ?
\begin{qoptions}
\opt{A}{If $n=8$, then $(\sqrt2-1)r<R$}
\opt{B}{If $n=12$, then $\sqrt2(\sqrt3+1)r>R$}
\opt{C}{If $n=4$, then $(\sqrt2-1)r<R$}
\opt{D}{If $n=5$, then $r<R$}
\end{qoptions}

\sect{Section-III: Single Digit Integer (0-9)}

\qn{1} Let $\alpha$ and $\beta$ be real numbers such that $-\dfrac{\pi}{4}<\beta<0<\alpha<\dfrac{\pi}{4}$. If $\sin(\alpha+\beta)=\dfrac14$ and $\cos(\alpha-\beta)=\dfrac34$, then the greatest integer less than or equal to
\[\left(\dfrac{\sin\alpha}{\cos\beta}+\dfrac{\cos\beta}{\sin\alpha}+\dfrac{\cos\alpha}{\sin\beta}+\dfrac{\sin\beta}{\cos\alpha}\right)^2\]
is\underline{\hspace{2em}}.

\qn{2} If $y(x)$ is the solution of the differential equation
\[xdy-(y^2-4y)dx=0\ \text{for}\ x>0,\ y(1)=2,\]
and the slope of the curve $y=y(x)$ is never zero, then the value of 17$y(2)$ is\underline{\hspace{2em}}.

\qn{3} The greatest integer less than or equal to
\[\int_1^4 \log_2(x^3+1)dx + \int_1^{\log_2 65}(2^x-1)^{1/3}dx\]
is $\lambda$ then sum of digits of $\lambda$ is\underline{\hspace{2em}}.

\qn{4} The product of all positive real values of x satisfying the equation
\[x^{\left(16(\log_2 x)^3-68\log_2 x\right)} = 2^{-16}\]
is\underline{\hspace{2em}}.

\qn{5} If
\[\beta = \lim_{x\to0} \dfrac{e^{x^3}-(1-x^3)^{1/3}+\left((1-x^2)^{1/2}-1\right)\tan x}{x\tan^2 x}\]
then the value of $6\beta$ is\underline{\hspace{2em}}.

\qn{6} Let $\beta$ be an integer. Consider the matrix
\[A = \begin{pmatrix}\beta & 0 & 1\\ 2 & 1 & -2\\ 3 & 1 & -2\end{pmatrix}\]
If $A^7-(\beta+1)A^6+\beta A^5$ is a singular matrix, then the value of $\beta$ is\underline{\hspace{2em}}.

\qn{7} Consider the hyperbola
\[\dfrac{x^2}{25}-\dfrac{y^2}{16}=1\]
with foci at S and S$_1$, where S lies on the positive x-axis. Let P be a point on the hyperbola, in the first quadrant. Let $\angle$SPS$_1=\alpha$, with $\alpha<\dfrac{\pi}{2}$. The straight line passing through the point S and having the same slope as that of the tangent at P to the hyperbola, intersects the straight line S$_1$P at P$_1$. Let $\delta$ be the distance of P from the straight line SP$_1$, and $\beta=$ S$_1$P. Then the greatest integer less than or equal to $\dfrac{\beta\delta}{9}\sin\dfrac{\alpha}{2}$ is\underline{\hspace{2em}}.

\qn{8} Consider the functions $f,g:\mathbb R\to\mathbb R$ defined by
\[f(x) = x^2+\dfrac{5}{12}\ \text{and}\ g(x) = \begin{cases}2\left(1-\dfrac{4|x|}{3}\right), & |x|\le\dfrac34,\\[4pt] 0, & |x|>\dfrac34.\end{cases}\]
If $\alpha$ is the area of the region
\[\left\{(x,y)\in\mathbb R\times\mathbb R : |x|\le\dfrac34,\ 0\le y\le\max\{f(x),g(x)\}\right\}\]
then the unit digit of $96\alpha$ is\underline{\hspace{2em}}.

\end{multicols}
\end{document}
